\chapter{Discusión y conclusiones finales}

\paragraph*{}Este último capítulo recoge una visión global del trabajo realizado. Se valora el grado de cumplimiento de los objetivos planteados inicialmente, las dificultades encontradas durante el desarrollo y las posibles líneas de mejora que podrían abordarse en el futuro. La verdad es que este cierre resulta especialmente importante en un proyecto de ampliación, porque permite distinguir con más claridad qué aportaba ya el sistema original y qué se ha incorporado para hacerlo más completo, flexible y útil.

\paragraph*{}El trabajo desarrollado ha partido de una aplicación previa orientada al análisis de supervivencia aplicado al abandono académico universitario. Sobre esa base se han añadido nuevos modelos, nuevas visualizaciones, soporte multilingüe, exportación de informes, análisis de variables no binarias y una integración más cuidada con la interpretación mediante inteligencia artificial. En conjunto, la aplicación resultante no sustituye el dashboard anterior, sino que lo amplía de forma progresiva y coherente con los objetivos definidos en el anteproyecto.

\section{Análisis de los objetivos iniciales}

\paragraph*{}El objetivo general era ampliar y mejorar una aplicación interactiva para el análisis de supervivencia aplicado al abandono académico universitario. Ese objetivo puede considerarse cumplido, aunque conviene ir más allá de la afirmación genérica y revisar qué ha aportado cada pieza.

\paragraph*{}El módulo de inteligencia artificial es probablemente la mejora más visible para el usuario. La aplicación permite solicitar interpretaciones asociadas a los resultados y presentarlas dentro de la propia interfaz, lo que añade una capa explicativa muy útil para quienes no dominan todos los detalles estadísticos de los modelos. Aun así, se ha mantenido un criterio claro desde el principio: la IA apoya la interpretación, no la sustituye. El análisis estadístico sigue siendo el núcleo, y la IA lo acompaña.

\paragraph*{}El soporte multilingüe en castellano e inglés puede parecer un detalle menor, pero en la práctica mejora la accesibilidad de la herramienta de forma real. Una aplicación que solo funciona en un idioma tiene un alcance limitado. Con este cambio, el dashboard puede usarse en contextos académicos más amplios sin fricciones innecesarias.

\paragraph*{}Ampliar el conjunto de variables fue otro objetivo relevante, y probablemente uno de los que más han enriquecido el análisis. Se han incorporado covariables como el género, la discapacidad, el tramo de edad, el nivel educativo previo y los créditos cursados. Esta última, \texttt{studied\_credits}, merece mención especial: no es binaria, lo que obligó a aplicar estrategias de agrupación y transformación según la técnica utilizada. Ese esfuerzo adicional se nota en los resultados, que ahora reflejan una realidad más matizada.

\paragraph*{}En cuanto a las técnicas estadísticas, la incorporación de los modelos Exponencial y Weibull junto con Random Survival Forest completa el panorama de forma coherente. El usuario puede ahora comparar una aproximación simple, un modelo paramétrico flexible y un enfoque basado en aprendizaje automático, todo dentro de la misma herramienta. Esa capacidad de comparación es, en sí misma, uno de los valores añadidos más claros del proyecto.

\paragraph*{}Por último, la exportación de informes PDF cierra el ciclo de forma práctica. No se trata solo de que los resultados sean visibles en pantalla, sino de que puedan llevarse fuera de la aplicación, compartirse, documentarse y usarse en contextos académicos o de gestión. Las pruebas realizadas confirman que la exportación funciona de forma controlada, evitando generar documentos vacíos cuando no se cumplen las condiciones necesarias.

\section{Problemas encontrados en el desarrollo}

\paragraph*{}No todo ha ido sobre ruedas, y sería poco honesto no reconocerlo. Algunos problemas eran esperables; otros aparecieron donde menos se anticipaban.

\paragraph*{}El primero y más transversal fue el de ampliar sin romper. Trabajar sobre una aplicación ya existente implica respetar lo que había: los callbacks de Dash, las páginas, el estado del dataset, las salidas ya definidas. Cada nueva funcionalidad debía encajar en esa estructura sin desestabilizarla. Eso obligó a revisar con calma la arquitectura del sistema antes de tocar nada, y a mantener una separación clara entre interfaz, análisis, visualización, exportación e interpretación.

\paragraph*{}La preparación de los datos también tuvo más peso del previsto. Las nuevas técnicas no podían aplicarse directamente sobre cualquier tabla: antes era necesario validar columnas, transformar tipos, eliminar registros incompletos y gestionar los casos en los que no había información suficiente para ajustar un modelo. Y cada variable tenía su particularidad. \texttt{studied\_credits} fue la más exigente, precisamente porque rompía el esquema binario en el que estaban pensadas muchas partes del sistema original.

\paragraph*{}Integrar librerías con requisitos distintos fue otro punto de fricción. \texttt{lifelines} para los modelos paramétricos y \texttt{scikit-survival} para Random Survival Forest no siempre se comportan igual ni devuelven resultados en el mismo formato. En el caso de RSF, construir la matriz de características y trabajar con estructuras específicas para datos censurados requirió más tiempo del esperado. La regresión de Cox también dio algún susto puntual con problemas de convergencia según las covariables seleccionadas, algo que se resolvió incorporando regularización y mecanismos de control.

\paragraph*{}La exportación PDF fue más compleja de lo que parecía a primera vista. Generar documentos con tablas, gráficos, textos e interpretaciones de forma estable no es trivial. Combinar figuras de Plotly con \texttt{Kaleido} y \texttt{ReportLab}, controlar tamaños y orientación de páginas, gestionar el contenido seleccionado por el usuario: cada detalle sumaba. Y el sistema tenía que comportarse bien tanto con resultados completos como con salidas parciales.

\paragraph*{}Por último, la integración de la IA introdujo una dependencia práctica inevitable. El modelo Qwen se ejecuta de forma local mediante \texttt{llama-server}, lo que tiene la ventaja de no depender de servicios externos, pero también implica que el servidor debe estar activo antes de solicitar una explicación. Gestionar esos errores de conexión, tiempos de espera y respuestas vacías sin que la aplicación se bloqueara fue una parte del desarrollo que no estaba del todo prevista al inicio.

\section{Conclusiones y posibles mejoras}

\paragraph*{}Mirando el resultado final, lo que más satisface no es la lista de funcionalidades añadidas, sino el hecho de que el proyecto tiene sentido como conjunto. La ampliación es coherente con el sistema original, las nuevas piezas encajan con las antiguas y el resultado es una herramienta más completa sin haber perdido lo que ya funcionaba.

\paragraph*{}Dicho eso, quedan cosas por hacer. Siempre quedan.

\paragraph*{}Una línea natural de mejora sería profundizar en el conjunto de variables analizadas, incorporando información sobre actividad en la plataforma, intentos previos, módulos cursados o indicadores socioeconómicos. Están en los datos, y podrían aportar mucho si se trabajan con cuidado.

\paragraph*{}También tendría sentido añadir validaciones más avanzadas de los modelos: métricas comparativas adicionales, validación cruzada, análisis de estabilidad. Eso permitiría valorar con más rigor qué técnica funciona mejor según el tipo de datos y el objetivo del usuario, en lugar de dejar esa decisión completamente en manos de quien usa la herramienta.

\paragraph*{}La interpretación mediante IA tiene también recorrido por delante. En una futura versión podría ser más contextualizada: no solo explicar qué resultado se ha obtenido, sino advertir sobre las precauciones que hay que tener al interpretarlo según el modelo y las variables elegidas. Una IA que no solo responde, sino que también avisa.

\paragraph*{}Y por último, el despliegue. Ahora mismo el sistema funciona en local, lo que limita su alcance real. Evolucionar hacia una aplicación en servidor, con gestión de usuarios, persistencia de análisis y almacenamiento de informes, abriría la puerta a un uso institucional genuino. Ese salto no es trivial, pero tampoco está tan lejos.

\paragraph*{}En definitiva, este trabajo ha conseguido lo que se proponía: acercar modelos estadísticos complejos a una interfaz comprensible y útil para explorar un problema educativo real. Eso es lo que queda cuando se apaga el ordenador, no las líneas de código ni las librerías instaladas, sino una herramienta que alguien podría usar de verdad para entender mejor por qué los estudiantes abandonan y, quizás, para hacer algo al respecto.